\chapter{漏洞防控系统}\label{chap:system}

本章给出第\ref{chap:method}章技术方案的系统化落地形态，重点说明模块划分、数据流与运行流程，使后续实现与复现实验具备清晰的工程参照。

\section{系统架构}

系统采用“数据准备—推理生成—评价筛选—发布回滚”的分层架构，可抽象为三类平面：
\begin{itemize}
	\item \textbf{数据平面}：负责多源数据采集、分块、去重、向量化与索引维护，提供可追溯的上下文检索能力。
	\item \textbf{控制平面}：负责任务编排与流程控制（RAG→CoT→Feedback→Evaluation），统一 schema 校验、重试策略、日志与结果落盘。
	\item \textbf{发布平面}：对通过门槛的策略进行发布、验证、灰度与回滚，并与集群变更流程（CI/CD 或运维流程）对接。
\end{itemize}

\section{核心模块}

\subsection{证据与知识库服务}
该模块面向“证据可追溯”目标，完成应用仓库与安全材料的收集、分块、索引与元数据维护，输出带来源标注的证据集合 $c=\{(t_i,\text{src}_i)\}$，并提供检索接口支持 RAG。

\subsection{策略生成与优化服务}
该模块以漏洞实例 $v$ 与上下文 $c$ 为输入，执行分阶段推导并输出候选策略集合 $\Pi_v$：
\begin{itemize}
	\item \textbf{RAG}：检索并汇聚与漏洞相关的证据片段，进行裁剪与去噪，保留来源标识。
	\item \textbf{CoT}：先输出结构化漏洞分析（类型、攻击向量、影响范围、风险类别），再在 \texttt{layer/method/actions} 约束下生成动作序列。
	\item \textbf{Feedback}：对候选进行评审，依据“缺失信息—前置条件—副作用—回滚边界”反馈进行二次优化。
\end{itemize}
同时集成 JSON/schema 合法性校验，确保产物可用于后续自动化处理。

\subsection{质量评价与审阅}
该模块实现门槛评分 $s(\pi)$ 与三维质量向量 $q(\pi)=(E,I,D)$ 的规则化检查，并在三个维度上生成排名 $R(\Pi_v,\text{dim})$。评价结论同时输出理由、证据来源与缺失信息清单，支持 human-in-the-loop 复核与决策。

\subsection{策略发布与回滚}
该模块将最终策略映射到可执行的控制手段（如准入策略、NetworkPolicy、运行时检测规则等），并提供“验证—灰度—回滚”的发布流程：上线前对配置语法、目标资源范围与冲突策略进行检查；上线后结合告警与观测数据评估副作用，必要时按预设边界回滚。

\section{工作流程}

端到端流程如下：
\begin{enumerate}
	\item 输入漏洞实例 $v$（CVE 描述与元数据），触发上下文检索构建 $c$。
	\item 基于 $v$ 与 $c$ 并行生成多份候选策略形成 $\Pi_v$，并进行 schema 校验与规范化。
	\item 对 $\Pi_v$ 执行门槛评分与三维评价，过滤未达标候选并对达标候选进行维度排名。
	\item 评价报告作为反馈约束触发候选优化；达标策略进入人工复核与发布决策。
	\item 发布到集群并执行验证与观测；若出现副作用或误报，按回滚边界撤销或调整策略。
\end{enumerate}

\section{部署与运行}

系统运行依赖包括：向量索引存储（FAISS 及元数据文件）、大模型调用接口与密钥管理、以及与集群交互的凭据与权限控制。为保证可运维性，建议记录每次生成的输入、检索证据摘要、候选策略、评分与排名结果，并对关键失败点（检索为空、schema 不合法、部署失败、回滚触发）进行可观测化告警。

